\ifdefined\gkver\else\def\gkver{student}\fi
\documentclass[\gkver]{gaokaozhenti}
\begin{document}

\chapter{2012年广东理综}
%% paperType: 理综物理部分

\begin{enumerate}
\item
%% number: 13
%% paperName: 2012年广东理综第 13 题 4 分
%% typeId: 1
%% type: 单选题
%% chapter: 气体性质
%% point: 分子力
%% method: 无
%% score: 4
%% degree: 650
%% duplicateId: 0
%% body:
清晨，草叶上的露珠是由空气中的水汽凝结成水珠，这一物理过程中，水分子间的\xzanswer{D}
\fourchoices[answer=D]
{引力消失，斥力增大}
{斥力消失，引力增大}
{引力、斥力都减小}
{引力、斥力都增大}
%% answer: D
%% memo:
\memoanswer{
水汽凝结成水珠，分子间距减小，分子引力和斥力都随分子间距的减小而增大，所以 D 选项正确，A、B、C 选项都是错误的。
}

\item
%% number: 14
%% paperName: 2012年广东理综第 14 题 4 分
%% typeId: 1
%% type: 单选题
%% chapter: 气体性质
%% point: 热力学第一定律
%% method: 无
%% score: 4
%% degree: 650
%% duplicateId: 0
%% body:
景颇族的祖先发明的点火器如图所示，用牛角做套筒，木质推杆前端粘着艾绒。猛推推杆，艾绒即可点燃，对筒内封闭的气体，再次压缩过程中\xzanswer{B}
\onepicture[h=1.2cm, align=right]{figs/14.png}
\fourchoices[answer=B]
{气体温度升高，压强不变}
{气体温度升高，压强变大}
{气体对外界做正功，其体内能增加}
{外界对气体做正功，气体内能减少}
%% answer: B
%% memo:
\memoanswer{
由于套筒内封闭着一定质量的气体，当猛推推杆时推杆迅速压缩气体，外界对气体做正功。由于这一过程进行得很快，可以看成是一个近似的绝热过程，即整个系统来不及向外界传递热量。根据热力学第一定律 $\Delta U=W+Q$，这时外力做的功只能用来增加气体的内能。这就使气体分子的运动加剧，引起气体分子平均动能增加，气体温度升高。所以艾绒即刻被点燃。由于被封闭的气体质量不变，温度升高，而体积变小，则由气体状态方程知压强变大。故选项 B 正确。
}

\item
%% number: 15
%% paperName: 2012年广东理综第 15 题 4 分
%% typeId: 1
%% type: 单选题
%% chapter: 磁场
%% point: 带电粒子在磁场中的运动
%% method: 无
%% score: 4
%% degree: 650
%% duplicateId: 0
%% body:
质量和电量都相等的带电粒子 M 和 N，以不同的速度率经小孔 S 垂直进入匀强磁场，运行的半圆轨迹如图中虚线所示，下列表述正确的是\xzanswer{A}
\onepicture[h=3.2cm, align=right]{figs/15.png}
\fourchoices[answer=A]
{M 带负电，N 带正电}
{M 的速度率小于 N 的速率}
{洛仑磁力对 M、N 做正功}
{M 的运行时间大于 N 的运行时间}
%% answer: A
%% memo:
\memoanswer{
由左手定则可知 M 带负电，N 带正电，选项 A 正确。由 $qvB=m\frac{v^{2}}{R}$ 得 $R=\frac{mv}{qB}$，由于二个带电粒子的质量和电量都相等，又进入到同一个匀强磁场中，由图及 A 选项的判断可知 $R_{N}<R_{M}$，故 $v_{N}<v_{M}$，所以 B 错误。由于洛伦兹力的方向始终与带电粒子的运动方向垂直，故洛伦兹力永远不会对 M、N 做功，则 C 选项错误。由 $T=\frac{2\pi R}{v}=\frac{2\pi m}{qB}$ 及题给条件可知，这二个带电粒子在磁场中运动的周期相等，又由图可见二个粒子在磁场中的偏转角相等，均偏转了半个周期，故在磁场中运动的时间相等，所以 D 错误。
}

\item
%% number: 16
%% paperName: 2012年广东理综第 16 题 4 分
%% typeId: 1
%% type: 单选题
%% chapter: 力物体的平衡
%% point: 三力平衡
%% method: 无
%% score: 4
%% degree: 650
%% duplicateId: 0
%% body:
如图所示，两根等长的轻绳将日光灯悬挂在天花板上，两绳与竖直方向的夹角为 $45^\circ$，日光灯保持水平，所受重力为$G$，左右两绳的拉力大小分别为\xzanswer{B}
\onepicture[h=2.4cm, align=right]{figs/16.png}
\fourchoices[answer=B]
{$G$和$G$}
{$\frac{\sqrt 2}{2}G$和$\frac{\sqrt 2}{2}G$}
{$\frac{1}{2}G$和$\frac{\sqrt 3}{2}G$}
{$\frac{1}{2}G$和$\frac{1}{2}G$}
%% answer: B
%% memo:
\memoanswer{
日光灯受重力 $G$ 及二根绳子拉力 $T$ 作用而处于平衡状态，则有 $G=2T\cos 45^\circ$，故 $T=\frac{\sqrt{2}}{2}G$，B 选项正确。
}

\item
%% number: 17
%% paperName: 2012年广东理综第 17 题 6 分
%% typeId: 2
%% type: 多选题
%% chapter: 曲线运动
%% point: 竖直平面圆周运动
%% method: 无
%% score: 6
%% degree: 450
%% duplicateId: 0
%% body:
如图是滑道压力测试的示意图，光滑圆弧轨道与光滑斜面相切，滑道底部$B$处安装一个压力传感器，其示数$N$表示该处所受压力的大小。某滑块从斜面上不同高度$h$处由静止下滑，通过$B$时，下列表述正确的有\xzanswer{BC}
\onepicture[h=3.0cm, align=right]{figs/17.png}
\fourchoices[answer=BC]
{$N$小于滑块重力}
{$N$大于滑块重力}
{$N$越大表明$h$越大}
{$N$越大表明$h$越小}
%% answer: BC
%% memo:
\memoanswer{
由动能定理 $mgh=\frac{1}{2}mv^{2}$，物体由斜面上 $h$ 高度处下滑到达 $B$ 处时的速度 $v=\sqrt{2gh}$，从 $B$ 处进入圆弧轨道后做圆周运动，在 $B$ 处，由牛顿第二定律得 $N-mg=m\frac{v^{2}}{R}$，故 $N=mg+m\frac{v^{2}}{R}=mg\left(1+\frac{2h}{R}\right)$，所以支持力 $N>G$，且 $h$ 越大 $N$ 越大，再由牛顿第三定律可知 B、C 正确。
}

\item
%% number: 18
%% paperName: 2012年广东理综第 18 题 6 分
%% typeId: 2
%% type: 多选题
%% chapter: 原子和原子核
%% point: 核裂变
%% method: 无
%% score: 6
%% degree: 650
%% duplicateId: 0
%% body:
能源是社会发展的基础，发展核能是解决能源问题的途径之一。下列释放核能的反应方程，表述正确的有\xzanswer{AC}
\fourchoices[answer=AC]
{\ce{^{3}_{1}H + ^{2}_{1}H -> ^{4}_{2}He + ^{1}_{0}n}是核聚变反应}
{\ce{^{3}_{1}H + ^{2}_{1}H -> ^{4}_{2}He + ^{1}_{0}n}是$\beta$衰变}
{\ce{^{235}_{92}U + ^{1}_{0}n -> ^{144}_{56}Ba + ^{89}_{36}Kr + 3^{1}_{0}n}是核裂变反应}
{\ce{^{235}_{92}U + ^{1}_{0}n -> ^{140}_{54}Xe + ^{94}_{38}Sr + 2^{1}_{0}n}是$\alpha$衰变}
%% answer: AC
%% memo:
\memoanswer{
两个轻核聚合成一个较重的核的反应为核聚变反应，故 A 对 B 错。重核被中子轰击后分裂成二个中等质量的原子核并放出若干个中子的反应为核裂变反应，故 C 对 D 错。原子核放出 $\alpha$ 或 $\beta$ 粒子后变成另一种原子核的反应称为原子核的衰变。原子核的衰变的反应物只有一种，故 A、C 选项正确。
}

\item
%% number: 19
%% paperName: 2012年广东理综第 19 题 6 分
%% typeId: 2
%% type: 多选题
%% chapter: 交流电
%% point: 交流电
%% method: 无
%% score: 6
%% degree: 830
%% duplicateId: 0
%% body:
某小型发电机产生的交变电动势为 $e=50\sin100\pi t$（V）。对此电动势，下列表述正确的有\xzanswer{CD}
\fourchoices[answer=CD]
{最大值是$50\sqrt 2\UV$}
{频率是$100\UHz$}
{有效值是$25\sqrt 2\UV$}
{周期是$0.02\Us$}
%% answer: CD
%% memo:
\memoanswer{
交变电动势瞬时值的表达式为 $e=E_{m}\sin \omega t=NBS\omega \sin \omega t=50\sin 100\pi t$，故交变电动势的最大值为 $50\UV$，有效值为 $E=\frac{E_{m}}{\sqrt{2}}=25\sqrt{2}\UV$，故 A 错 C 对。由 $\omega=2\pi f=100\pi$ 得频率是 $f=50\UHz$，周期 $T=\frac{1}{f}=0.02\Us$，故 B 错 D 对。
}

\item
%% number: 20
%% paperName: 2012年广东理综第 20 题 6 分
%% typeId: 2
%% type: 多选题
%% chapter: 电场
%% point: 带电粒子的偏转
%% method: 无
%% score: 6
%% degree: 650
%% duplicateId: 0
%% body:
如图是某种静电矿料分选器的原理示意图，带电矿粉经漏斗落入水平匀强电场后，分落在收集板中央的两侧，对矿粉分离的过程，下列表述正确的有\xzanswer{BD}
\onepicture[h=3.2cm, align=right]{figs/20.png}
\fourchoices[answer=BD]
{带正电的矿粉落在右侧}
{电场力对矿粉做正功}
{带负电的矿粉电势能变大}
{带正电的矿粉电势能变小}
%% answer: BD
%% memo:
\memoanswer{
正、负电的矿粉受电场力做正功分别落到左、右两侧，根据电场力做功与电势能的变化关系：电场力做正功，电势能减小。故正、负电的矿粉的电势能都变小，B、D 正确。
}

\item
%% number: 21
%% paperName: 2012年广东理综第 21 题 6 分
%% typeId: 2
%% type: 多选题
%% chapter: 曲线运动
%% point: 人造地球卫星
%% method: 无
%% score: 6
%% degree: 450
%% duplicateId: 0
%% body:
如图所示，飞船从轨道1变轨至轨道2。若飞船在两轨道上都做匀速圆周运动，不考虑质量变化，相对于在轨道1上，飞船在轨道2上的\xzanswer{CD}
\onepicture[h=3.2cm, align=right]{figs/21.png}
\fourchoices[answer=CD]
{动能大}
{向心加速度大}
{运行周期长}
{角速度小}
%% answer: CD
%% memo:
\memoanswer{
由万有引力定律及向心力公式得 $G\frac{mM}{r^{2}}=ma=m\frac{v^{2}}{r}=mr\omega^{2}=mr\frac{4\pi^{2}}{T^{2}}$，由题知 $r_{2}>r_{1}$，由此可知 $E_{K}=\frac{1}{2}mv^{2}=\frac{GmM}{2r}$，则 $E_{k2}<E_{k1}$，A 错。$a=\frac{GM}{r^{2}}$，则 $a_{2}<a_{1}$，B 错。$\omega=\sqrt{\frac{GM}{r^{3}}}$，则 $\omega_{2}<\omega_{1}$，D 正确。$T=\frac{2\pi}{\omega}$，则 $T_{2}>T_{1}$，C 正确。
}

\item
%% number: 34
%% paperName: 2012年广东理综第 34 题 8 分
%% typeId: 4
%% type: 实验题
%% chapter: 电路
%% point: 伏安法测电阻
%% method: 无
%% score: 8
%% degree: 650
%% duplicateId: 0
%% body:
某同学测量一个圆柱体的电阻率，需要测量圆柱体的尺寸和电阻。
\begin{enumerate}
	\item
	分别使用游标卡尺和螺旋测微器测量圆柱体的长度和直径，某次测量的示数如图（a）和图（b）所示，长度为\_\_\_\_\_\_\_cm，直径为\_\_\_\_\_\_\_mm。
	\item
	按图（$c$）连接电路后，实验操作如下：

	（$a$）将滑动变阻器$R_{1}$的阻值置于最\_\_\_\_\_处（填“大”或“小”）；将$S_{2}$拨向接点1，闭合$S_{1}$，调节$R_{1}$，使电流表示数为$I_{0}$；

	（$b$）将电阻箱$R_{2}$的阻值调至最\_\_\_\_\_\_（填“大”或“小”）；将$S_{2}$拨向接点2；保持$R_{1}$不变，调节$R_{2}$，使电流表示数仍为$I_{0}$，此时$R_{2}$阻值为$1280\,\UO$；
	\item
	由此可知，圆柱体的电阻为\_\_\_\_\_\_$\UO$。
\end{enumerate}
\threepicture[align=right, hA=1.15cm, hB=1.5cm, hC=1.75cm]
{figs/34_1a.png}
{figs/34_1b.png}
{figs/34_1c.png}
%% answer: 
\jdanswer{
\begin{enumerate}
	\item
	$5.02$，$5.315$

	\item
	大，大

	\item
	$1280$

\end{enumerate}
}
%% memo:
\memoanswer{
①5.02，5.315

②电学实验的设计要遵循科学性原则、安全性原则和准确性原则。此电路中滑动变阻器是以限流方式接入电路中的，故在（a）步骤中合上开关前应使其接入电路中的阻值为最大，以保证电路安全。同理（b）步骤中亦将电阻箱 $R_{2}$ 的阻值调至最大。③步骤中，由闭合电路欧姆定律得 $I_{0}=\frac{E}{R+R_{1}+R_{g}+r}$，其中 $R$ 表示圆柱体的电阻。步骤中，仍由闭合电路欧姆定律得 $I_{0}=\frac{E}{R_{2}+R_{1}+R_{g}+r}$，由等量代换可得 $R=R_{2}=1280\UO$。
}

\item
%% number: 34
%% paperName: 2012年广东理综第 34 题 10 分
%% typeId: 4
%% type: 实验题
%% chapter: 力物体的平衡
%% point: 胡克定律
%% method: 无
%% score: 10
%% degree: 650
%% duplicateId: 0
%% body:
某同学探究弹力与弹簧伸长量的关系。
\begin{enumerate}
	\item
	将弹簧悬挂在铁架台上，将刻度尺固定在弹簧一侧，弹簧轴线和刻度尺都应在\_\_\_\_\_\_方向（填“水平”或“竖直”）。
	\item
	弹簧自然悬挂，待弹簧\_\_\_\_\_\_时，长度记为$L_{0}$，弹簧下端挂上砝码盘时，长度记为$L_{x}$；在砝码盘中每次增加 $10\Ug$ 砝码，弹簧长度依次记为$L_{1}$至$L_{6}$，数据如下表：
	\begin{center}
	\begin{tabular}{|c|c|c|c|c|c|c|c|c|}
	\hline
	代表符号 & $L_{0}$ & $L_{x}$ & $L_{1}$ & $L_{2}$ & $L_{3}$ & $L_{4}$ & $L_{5}$ & $L_{6}$ \\ \hline
	数值（cm） & 25.35 & 27.35 & 29.35 & 31.30 & 33.4 & 35.35 & 37.40 & 39.30 \\ \hline
	\end{tabular}
	\end{center}

	表中有一个数值记录不规范，代表符号为\_\_\_\_\_\_\_。由表可知所用刻度尺的最小长度为\_\_\_\_\_\_。
	\item
	如图是该同学根据表中数据作的图，纵轴是砝码的质量，横轴是弹簧长度与\_\_\_\_\_\_\_\_\_的差值（填“$L_{0}$”或“$L_{1}$”）。
	\item
	由图可知弹簧的劲度系数为\_\_\_\_\_\_\_\_\_$\UNm$；通过图和表可知砝码盘的质量为\_\_\_\_\_\_\_\_\_$\Ug$（结果保留两位有效数字，重力加速度取 $9.8\Umsq$）。
\end{enumerate}
\onepicture[h=3.2cm, align=right]{figs/34_2.png}
%% answer: 
\jdanswer{
\begin{enumerate}
	\item
	竖直

	\item
	稳定，$L_{3}$，$1\Umm$

	\item
	$L_{0}$

	\item
	$4.9$，$10$

\end{enumerate}
}
%% memo:
\memoanswer{
①竖直。②测量仪器是10分度的需要估读到最小精度的下一位，则此刻度尺的最小精度应为 $1\Umm$，故 $L_{3}$ 读数错误。③横轴表示的是弹簧的形变量，故为弹簧长度与 $L_{0}$ 的差值。④由 $mg=kx$ 得图像的斜率 $k'=\frac{m}{x}=\frac{k}{g}=\frac{10\times 10^{-3}}{2\times 10^{-2}}=0.5$，故 $k=4.9\UNm$，挂了砝码盘后弹簧伸长了 $2\Ucm$，由 $mg=kx$ 知其质量为 $10\Ug$。
}

\item
%% number: 35
%% paperName: 2012年广东理综第 35 题 18 分
%% typeId: 6
%% type: 计算题
%% chapter: 电磁感应
%% point: 电磁感应电路分析
%% method: 无
%% score: 18
%% degree: 450
%% duplicateId: 0
%% body:
如图所示，质量为$M$的导体棒$ab$，垂直放在相距为$l$的平行光滑金属导轨上。导轨平面与水平面的夹角为$\theta$，并处于磁感应强度大小为$B$、方向垂直与导轨平面向上的匀强磁场中。左侧是水平放置、间距为$d$的平行金属板。$R$和$R_{x}$分别表示定值电阻和滑动变阻器的阻值，不计其他电阻。
\begin{enumerate}
	\item
	调节$R_{x}=R$，释放导体棒，当棒沿导轨匀速下滑时，求通过棒的电流$I$及棒的速率$v$。
	\item
	改变$R_{x}$，待棒沿导轨再次匀速下滑后，将质量为$m$、带电量为$+q$的微粒水平射入金属板间，若它能匀速通过，求此时的$R_{x}$。
\end{enumerate}
\onepicture[h=3.2cm, align=right]{figs/35.png}
%% answer: 
\jdanswer{
\begin{enumerate}
	\item
	$I=\frac{Mg\sin\theta}{Bl}$，$v=\frac{2MgR\sin \theta }{B^{2}l^{2}}$

	\item
	$R_{x}=\frac{mldB}{Mq\sin\theta}$

\end{enumerate}
}
%% memo:
\memoanswer{
（1）当 $R_{x}=R$ 棒沿导轨匀速下滑时，由平衡条件 $Mg\sin\theta=F$

安培力 $F=BIl$

解得 $I=\frac{Mg\sin\theta}{Bl}$

感应电动势 $E=Blv$

电流 $I=\frac{E}{2R}$

解得 $v=\frac{2MgR\sin\theta}{B^{2}l^{2}}$

（2）微粒水平射入金属板间，能匀速通过，由平衡条件 $mg=q\frac{U}{d}$

棒沿导轨匀速，由平衡条件 $Mg\sin \theta=BI_{1}l$

金属板间电压 $U=I_{1}R_{x}$

解得 $R_{x}=\frac{mldB}{Mq\sin\theta}$
}

\item
%% number: 36
%% paperName: 2012年广东理综第 36 题 18 分
%% typeId: 6
%% type: 计算题
%% chapter: 功和能
%% point: 动能定理
%% method: 无
%% score: 18
%% degree: 250
%% duplicateId: 0
%% body:
如图（a）所示的装置中，小物块A、B质量均为$m$，水平面上PQ段长为$l$，与物块间的动摩擦因数为$\mu$，其余段光滑。初始时，挡板上的轻质弹簧处于原长；长为$r$的连杆位于图中虚线位置；A紧靠滑杆（A、B间距大于2$r$）。随后，连杆以角速度$\omega$匀速转动，带动滑杆作水平运动，滑杆的速度-时间图像如图（b）所示。$A$在滑杆推动下运动，并在脱离滑杆后与静止的B发生完全非弹性碰撞。
\begin{enumerate}
	\item
	求A脱离滑杆时的速度$u_{0}$，及A与B碰撞过程的机械能损失$\Delta E$。
	\item
	如果AB不能与弹簧相碰，设AB从$P$点到运动停止所用的时间为$t_{1}$，求$\omega$的取值范围，及$t_{1}$与$\omega$的关系式。
	\item
	如果AB能与弹簧相碰，但不能返回到P点左侧，设每次压缩弹簧过程中弹簧的最大弹性势能为$E_{p}$，求$\omega$的取值范围，及$E_{p}$与$\omega$的关系式（弹簧始终在弹性限度内）。
\end{enumerate}
\twopicture[align=right, hA=2.4cm, hB=2.4cm]
{figs/36a.png}
{figs/36b.png}
%% answer: 
\jdanswer{
\begin{enumerate}
	\item
	$u_{0}=\omega r$，$\Delta E=\frac{1}{8}m\omega^{2}r^{2}$

	\item
	$0<\omega\leq\frac{4l}{rt_{1}}$，$t_{1}=\frac{\omega r}{2\mu g}$

	\item
	$E_{p}=\frac{m(\omega^{2}r^{2}-8\mu gl)}{4}$

\end{enumerate}
}
%% memo:
\memoanswer{
（1）由题知，A 脱离滑杆时的速度 $u_{0}=\omega r$

设 A、B 碰后的速度为 $v_{1}$，由动量守恒定律

$mu_{0}=2mv_{1}$

A 与 B 碰撞过程损失的机械能 $\Delta E=\frac{1}{2}mu_{0}^{2}-\frac{1}{2}\times 2mv_{1}^{2}$

解得 $\Delta E=\frac{1}{8}m\omega^{2}r^{2}$

（2）AB 不能与弹簧相碰，设 AB 在 PQ 上运动的加速度大小为 $a$，由牛顿第二定律及运动学规律

$\mu\cdot 2mg=2ma$

$v_{1}=at_{1}$

$x=\frac{v_{1}}{2}t_{1}$

由题知 $x\leq l$

联立解得 $0<\omega\leq\frac{4l}{rt_{1}}$

$t_{1}=\frac{\omega r}{2\mu g}$

（3）AB 能与弹簧相碰 $\mu\cdot 2mgl<\frac{1}{2}\times 2mv_{1}^{2}$

不能返回到 P 点左侧 $\mu\cdot 2mg\cdot 2l\geq\frac{1}{2}\times 2mv_{1}^{2}$

解得 $\frac{2\sqrt{2\mu gl}}{r}<\omega\leq\frac{4\sqrt{\mu gl}}{r}$

AB 在 Q 点速度为 $v_{2}$，AB 碰后到达 Q 点过程，由动能定理

$-\mu\cdot 2mgl=\frac{1}{2}\times 2mv_{2}^{2}-\frac{1}{2}\times 2mv_{1}^{2}$

AB 与弹簧接触到压缩最短过程，由能量守恒

$E_{p}=\frac{1}{2}\times 2mv_{2}^{2}$

解得 $E_{p}=\frac{m(\omega^{2}r^{2}-8\mu gl)}{4}$
}

\end{enumerate}

\end{document}
